\label{ssub:branching-eval}
The corrected branching evaluation uses a temporal case split with an explicit
held-out partition. The older \texttt{final\_composite\_branching.pkl} artifact
was trained on the full log and is therefore treated as historical deployment
material, not as held-out evaluation evidence. The evaluation artifact used for
the optimized final branching candidate is
\texttt{composite\_branching\_evaluation\_train70\_rfopt\_v1.pkl}; it is trained
only on the outer development split and has zero case overlap with the held-out
test and with the 76 fixed-replay routes. The earlier
\texttt{composite\_branching\_evaluation\_train70.pkl} artifact is preserved as
the pre-optimization corrected baseline.

The split contains $22{,}056$ outer-development cases and $9{,}453$ held-out
cases. Inside the development partition, $18{,}747$ cases are used as inner
training cases and $3{,}309$ as validation cases. The fixed-replay case set is
fully contained in the held-out partition. The common BPMN-replay dataset is
constructed by replaying complete lifecycle events on
\texttt{models/v4\_replay.bpmn}. A row is evaluated only when the current event
can be synchronized and fired, the resulting marking exposes multiple enabled
BPMN candidates, and the observed next activity is one of those candidates.

\begin{table}[H]
    \caption{Corrected BPMN-replay branching dataset and split.}
    \centering
    \renewcommand{\arraystretch}{1.15}
    \scriptsize
    \begin{tabular}{lr}
        \toprule
        \textbf{Quantity} & \textbf{Value} \\
        \midrule
        Outer-development cases & $22{,}056$ \\
        Held-out cases & $9{,}453$ \\
        Inner-train / validation cases & $18{,}747$ / $3{,}309$ \\
        Fixed-replay held-out cases & $76$ \\
        Complete lifecycle events inspected & $475{,}306$ \\
        Synchronized events & $228{,}895$ \\
        Deterministic single-candidate states & $103{,}337$ \\
        Multi-candidate decision observations & $98{,}332$ \\
        Held-out evaluated observations & $29{,}385$ \\
        True label not enabled after replay & $24{,}526$ \\
        Case overlap with evaluation training & $0$ \\
        \bottomrule
    \end{tabular}
    \label{tab:branching_corrected_dataset_joao}
\end{table}

All branching methods in Table~\ref{tab:branching_corrected_methods_joao} are
evaluated on the same held-out BPMN-replay observations, candidate sets and
targets. The Random-Forest classifier is reported both as a raw classifier and
as the BPMN-constrained runtime policy. In this corrected evaluation, the raw
and constrained scores are equal because the RF top classes were already valid
under the candidate sets used by the held-out rows. The RF is not calibrated;
its probabilities are used as ranking scores, not as calibrated branch
probabilities.

\begin{table}[H]
    \caption{Corrected branching method comparison on the common held-out
    BPMN-replay dataset ($29{,}385$ observations).}
    \centering
    \renewcommand{\arraystretch}{1.1}
    \scriptsize
    \begin{tabular}{lrrrr}
        \toprule
        \textbf{Method} & \textbf{Accuracy} & \textbf{Macro-F1} &
        \textbf{Weighted-F1} & \textbf{Fallback / uncovered} \\
        \midrule
        Majority baseline & $0.838455$ & $0.609842$ & $0.822150$ & $0.000$ \\
        Random candidate baseline & $0.434848$ & $0.304639$ & $0.536529$ & $0.000$ \\
        Probability branching & $0.778526$ & $0.514706$ & $0.778638$ & $0.000$ \\
        RF raw classifier & $0.935341$ & $0.707269$ & $0.935807$ & $0.000$ \\
        RF BPMN-constrained & $0.935341$ & $0.707269$ & $0.935807$ & $0.000$ \\
        Composite runtime & $0.935341$ & $0.707269$ & $0.935807$ & $0.000$ \\
        \bottomrule
    \end{tabular}
    \label{tab:branching_corrected_methods_joao}
\end{table}

The optimized results change the interpretation of the advanced branching
claims. The predictive branching implementation is a genuine Advanced-II
component: it is trained without held-out case leakage, consumes case history
and runtime-equivalent features, and is constrained to BPMN-enabled candidates.
It does not support a claim of perfect prediction or statistical superiority.
The Random Forest improves macro-F1 over the majority baseline and is used as
the first layer of the composite runtime policy.

The Basic probabilistic component is methodologically clean after correction:
it uses training-only empirical distributions, BPMN candidate filtering,
state-conditioned support when available, and explicit backoff. In the final
scope, this component is reported as the Basic method and as the fallback layer
of the Composite runtime.

The transition-aware smoke test remains separate from the RF-composite
generative smoke. The original generative smoke validates BPMN transition
execution and resource-allocation integration with an empty composite hierarchy.
The new RF-composite smoke loads the leakage-free evaluation artifact and checks
that the RF composite can participate in a generative simulation. It is
integration evidence only and is not used as a statistical ranking benchmark.

Because the optimized RF artifact changes runtime branch choices, the fixed
replay was rerun for all five allocation methods on the same 76 held-out routes
and seeds. The rerun produced no failed runs. It also changed the ordering by
conditional mean cycle time: ParkSong-Composite and RoundRobin completed all
fixed routes in the optimized candidate, whereas the pre-optimization corrected
run censored one route for each of those strategies. Therefore, fixed-replay
timing differences must be interpreted jointly with completion and censorship
counts rather than as an isolated ranking claim.
